\documentclass{article}
\usepackage{graphicx} % Required for inserting images

\title{Sutton and Barton Exercises Chapter 4}
\author{danieltocco0 }
\date{October 2025}

\begin{document}

\maketitle

\section{Chapter 5}
\textit{Exercise 5.1} Consider the diagrams on the right in Figure 5.1. Why does the estimated
value function jump up for the last two rows in the rear? Why does it drop off for the
whole last row on the left? Why are the frontmost values higher in the upper diagrams
than in the lower?
\\
\\

Because the player sticks only on 20 or 21, whereas the dealer sticks on 17

or above. \\

As to the second question, the dealer has the highest probability of winning 

if they are showing an ace. \\

As to the third question, having a usable ace gives the player (and dealer, 

but moreso player due to the different policies), a higher probability of 

winning (to be more specific, the player sum is 12 in the front row and 

if he has a usable ace, he has more leeway to hit than if the ace is not usable). 
\\\\
\textit{Exercise 5.2} Suppose every-visit MC was used instead of first-visit MC on the blackjack
task. Would you expect the results to be very different? Why or why not?
\\
\\

No. The states are the dealer's show card and your cards. The dealer's 

show card is always constant, so averaging that does nothing. Your cards 

will never revisit the same state because there is no zero card.
\\\\
\textit{Exercise 5.3} What is the backup diagram for Monte Carlo estimation of $q_{\pi}$? 
\\
\\

Not sure how to draw that here in LaTex without spending more than it's 

worth, but it will look similar to the one state-value diagram. It will 

show the entire episode of actions for only the actions that are sampled 

with the first action as the root.
\\\\
\textit{Exercise 5.4} The pseudocode for Monte Carlo ES is inefficient because, for each state–
action pair, it maintains a list of all returns and repeatedly calculates their mean. It would
be more efficient to use techniques similar to those explained in Section 2.4 to maintain
just the mean and a count (for each state–action pair) and update them incrementally.
Describe how the pseudocode would be altered to achieve this.  
\\
\\
From section 2.4, equation 2.3, the incremental update rule:\\

$Q_{n+1} = \frac{1}{n}[R_{n} - Q_{n}]$ \\

The updated pseudocode then (all lines not shown are the same): \\

Append G to $\textit{Returns}(S_{t}, A_{t}) \Rightarrow \\
\indent N(S_{t}, A_{t}) \leftarrow N(S_{t}, A_{t}) + 1$
\\

$Q(S_{t}, A_{t}) \leftarrow average(\textit{Returns}(S_{t}, A_{t})) \Rightarrow\\
\indent Q(S_{t}, A_{t}) \leftarrow Q(S_{t}, A_{t}) + \frac{1}{N}[G_{n} - Q_{n}]$
\\\\
\textit{Exercise 5.5} Consider an MDP with a single nonterminal state and a single action
that transitions back to the nonterminal state with probability p and transitions to the
terminal state with probability 1 p . Let the reward be +1 on all transitions, and let
 = 1. Suppose you observe one episode that lasts 10 steps, with a return of 10. What
are the first-visit and every-visit estimators of the value of the nonterminal state?  
\\
\\

Originally, I was attempting to multiply by the probabilities. I missed the 

key insight of Monte Carlo Methods, in that we are learning from experience 

and the probabilities are built into the empirical observations. Thus, we have:

$G_{0} = 10$ for first-visit, \\

$\frac{1}{n}\sum_{n=0}^{n=9}G_{i} = \frac{(10 + 9 + 8 + 7 + 6 + 5 + 4 + 3 + 2 + 1)}{10} = 5.5$ for every-visit
\\\\
\textit{Exercise 5.6} What is the equation analogous to (5.6) for action values Q(s,a) instead of state values V(s), again given returns generated using b? 
\\
\\
$Q(s,a) \doteq \frac{\Sigma_{t \in \tau_{s,a}}\rho_{t:T(t)-1}G_{t}}{\tau(s,a)}$
\\\\
\textit{Exercise 5.7} In learning curves such as those shown in Figure 5.3 error generally decreases with training, as indeed happened for the ordinary importance-sampling method. But for the weighted importance-sampling method error first increased and then decreased. Why do you think this happened? 
\\
\\

Because initially you may have some extremely high ratios, 

and since it's weighted and you're near the start with few samples, 

those specific instances may dominate. In other words, the variance 

is high to start. Think of, for whatever reason you'd get a ratio 

of 100 vs 1 vs 2... the ratio weight for 100/(100+1+2) is going 

to dominate. As you increase the amount of samples, this

weight's value will decrease... And as discussed, you eventually converge. 
\\\\
\textit{Exercise 5.8} The results with Example 5.5 and shown in Figure 5.4 used a first-visit MC method. Suppose that instead an every-visit MC method was used on the same problem. Would the variance of the estimator still be infinite? Why or why not? 
\\
\\

Yes, you still have infinite variance. In this case, you 

take the sum of square of the sum of all ratios. E.g., if you had 

an episode of length 3, you'd get $(\rho_{1})^{2} + (\rho_{2}\rho_{1})^{2} +  (\rho_{3}\rho_{2}\rho_{1})^{2}$ (and yes, 

I got lazy; there are factors the probabilities that must go in front of these 

terms). In short, we at least recover the first-visit example sum, 

which diverges, so we know this in consequence diverges too.
\\\\
\textit{Exercise 5.9} Exercise 5.9 Modify the algorithm for first-visit MC policy evaluation (Section 5.1) to use the incremental implementation for sample averages described in Section 2.4 
\\
\\

This is similar to exercise 5.4: \\

Append G to $\textit{Returns}(S_{t}, A_{t}) \Rightarrow \\
\indent N(S_{t}) \leftarrow N(S_{t}) + 1$
\\

$V(S_{t}) \leftarrow average(\textit{Returns}(S_{t})) \Rightarrow\\
\indent V(S_{t}) \leftarrow V(S_{t}) + \frac{1}{N}[G_{n} - V(S_{t})]$

As before, all lines not shown are the same.
\\\\
\textit{Exercise 5.10}  Derive the weighted-average update rule (5.8) from(5.7). Follow the pattern of the derivation of the unweighted rule (2.3). 
\\
\\

$V_{n} \doteq \frac{\sum_{k=1}^{n-1} W_{k}G_{k}}{\sum_{k=1}^{n-1} W_{k}}, n \geq 2$ \indent \Rightarrow
\\\\

$V_{n+1} \doteq \frac{C_{n}V_{n} + W_{n+1}G_{n+1}}{C_{n} + W_{n+1}} \indent \Rightarrow
\\\\
\indent = \frac{C_{n}V_{n} + W_{n+1}G_{n+1}}{C_{n} + W_{n+1}} \indent \Rightarrow
\\\\
\indent V_{n+1} - V_{n} \doteq \frac{C_{n}V_{n} + W_{n+1}G_{n+1}}{C_{n} + W_{n+1}} - V_{n}
\\\\
\indent = \frac{C_{n}V_{n} + W_{n+1}G_{n+1} - V_{n}(C_{n} + W_{n+1})}{C_{n} + W_{n+1}}
\\\\
\indent = \frac{W_{n+1}(G_{n+1} - V_{n})}{C_{n} + W_{n+1}} \indent \Rightarrow
\\\\
\indent V_{n+1} = V_{n} + \frac{W_{n+1}(G_{n+1} - V_{n})}{C_{n} + W_{n+1}}
\\\\
\indent = V_{n} + \frac{W_{n+1}}{C_{n} + W_{n+1}}(G_{n+1} - V_{n})$
\\\\

This is the same as 5.8, except the indexing is different. I 

used GPT, of course, to help break this down for me; the 

$W_{n+1}$ is already absorbed into the $C_{n}$ term in 

5.8 (as is the $W_{n}$). 
\\\\
\textit{Exercise 5.11} In the boxed algorithm for off-policy MC control, you may have been expecting the W update to have involved the importance-sampling ratio $\frac{\pi(A_{t}|S_{t})}{b(A_{t}|S_{t})}$, but instead it involves $\frac{1}{b(A_{t}|S_{t})}$ . Why is this nevertheless correct?
\\
\\

Because in the preceding line, we have $A_{t} \neq \pi(S_{t})$, 

then exit the inner loop. This is saying that if the the action

followed by B is not the optimal action, then we exit the inner 

loop. The policy $\pi(S_{t})$ takes $argmax_{a}Q(S_{t}, a)$, which

has probability 1 (with ties broken consistently). Therefore, that

is how we end up with $\frac{1}{b(A_{t}|S_{t})}.$
\\\\
\textit{Exercise 5.12} Show the steps to derive (5.14) from (5.12).
\\
\\

Originally, I assumed E[XY] = E[X] + E[Y] applied, but the terms 

X and Y must be independent; we need a stronger assumption here 

which is the law of iterated expectation, in which we condition 

up to the history of t+k, giving:\\

$\mathbb{E_{b}}[\mathbb{E_{b}}[\rho_{t:T-1}R_{t+k}|F_{t+k}]]$ \indent = \\
\indent $\mathbb{E_{b}}[R_{r+k}\mathbb{E_{b}}[\rho_{t:T-1}|F_{t+k}]]$
\\

Now we can split the importance sampling ratios, and because we're given

the history \textit{up to} $F_{t+k}$, we can take the applicable ratio out: 

\\
\indent $\rho_{t:t+k-1}\rho_{t+k:T-1} \Rightarrow$ \\
\indent $\rho_{t:t+k-1}\mathbb{E_{b}[\rho_{t+k:T-1}|F_{t+k}] \Rightarrow$ \\
\indent $\mathbb{E_{b}}[R_{r+k}\rho_{t:t+k-1}\mathbb{E_{b}}[\rho_{t+k:T-1}|F_{t+k}]]$
\\\\

Now we can split the ratios and condition recursively as before:
\\
\indent $\rho_{t:t+k}\rho_{t+k+1:T-1} \indent \Rightarrow$ \\
\indent $\mathbb{E_{b}}[R_{t+k}\rho_{t:t+k-1}\mathbb{E_{b}[\rho_{t+k}\mathbb{E_{b}}[\rho_{t+k+1:T-1}|F_{t+k+1}]|F_{t+k}}}]]
\\

Now since $\mathbb{E_{b}}[\rho_{t+k}]$ is $\Sigma_{a}b(A_{t+k}|S_{t+k})$ $\frac{\pi(A_{t+k}|S_{t+k})}{b(A_{t+k}|S_{t+k})}$, we know this evaluates 

to 1, and knowing we can do this recursively for every term all the 

way up to the end, we know they all just evaluate to 1, so we 

end up with 5.14: \\

$\mathbb{E_{b}}[R_{t+k}\rho_{t:t+k-1}].



\\\\
\end{document}
